\section{Questions and scope}\label{sec:prelim}

Let $\X\in[0,1]^{3\times T\times H\times W}$ be a low-light clip and $\Y$ its
registered well-lit target. A frozen encoder $\Phi$ supplies token fields
$Z^{(\ell)}=\Phi^{(\ell)}(\X)$ to a trainable pixel branch
$\mathcal P_\theta$, which predicts $\Yhat$.
We ask three questions under a fixed reconstruction protocol.

\paragraph{Does conditioning help?}
We train the same pixel branch from scratch with either pretrained tokens or
zero tokens. The zero-token arm retains the conditioning modules and their
parameters. This controls the branch architecture and training schedule, while
testing the benefit of providing the pretrained representation. Attributing
that benefit specifically to pretraining also requires a correctly reproduced
random-encoder arm. We report its corrected single-seed evaluation separately
and distinguish perceptual from fidelity evidence (Sec.~\ref{sec:exp:control}).

\paragraph{Which parts of the field matter after training?}
We apply spatial averaging, temporal averaging, and principal-component
projection to a model trained with the full field. Each intervention is
compared with its own identity control on the same checkpoint and clips.
Preserved performance establishes tolerance to that particular replacement.
A degradation can reflect lost information or a shift in conditioning, so it
cannot alone identify the minimum representation needed for training.

\paragraph{Does temporal information improve the reconstruction?}
We first test whether the frozen encoder responds to frame order, then test
whether the restoration does. Finally, we add a loss that directly matches
temporal differences in the frozen feature space. Improving its diagnostic is
distinct from improving perceptual quality or pixel-space temporal behaviour.
These questions are evaluated at the final checkpoint of a fixed schedule;
short-budget experiments are exploratory unless explicitly identified.
